\begin{proposition}[Explicit witnesses for the two singular examples]
\label{prop:elliott-explicit-witnesses}
Put $f(X)=X^2+X$.
\begin{enumerate}
\item Let $D=\mathbb F_2[[t^2,t^3]]$. The polynomials
\begin{align*}
 h_2&=\frac{f^4+t^4f^2}{t^{10}},&
 h_1&=\frac{f}{t^2}+t h_2
\end{align*}
belong to $\operatorname{Int}(D)$ and satisfy
\begin{equation}\label{eq:cusp-explicit-witness}
 X^2+X=t^2h_1+t^3h_2.
\end{equation}
\item Let $D=\mathbb F_2+t\mathbb F_4[[t]]$, where
$\mathbb F_4=\mathbb F_2(\alpha)$ and $\alpha^2+\alpha+1=0$.
The polynomials
\begin{align*}
 k_2&=\frac{f^2+tf}{t^2},&
 k_1&=\frac{\alpha f^2+\alpha^2tf}{t^2}
\end{align*}
belong to $\operatorname{Int}(D)$ and satisfy
\begin{equation}\label{eq:residue-extension-explicit-witness}
 X^2+X=t k_1+\alpha t k_2.
\end{equation}
\end{enumerate}
In particular $X^2+X\in\mathfrak m\operatorname{Int}(D)$ in both
examples, where $\mathfrak m$ is the maximal ideal of $D$.
\end{proposition}
\begin{proof}
For the first ring, put $V=\mathbb F_2[[t]]$. Separating even and
odd exponents gives $D=\mathbb F_2+t^2V$, with maximal ideal
$(t^2,t^3)$. For $x=a+t^2y\in D$, where $a\in\mathbb F_2$ and
$y\in V$, put $g=f(x)/t^2=y+t^2y^2\in V$. Write
$g=b+ct+t^2z$ with $b,c\in\mathbb F_2$ and $z\in V$. Then
$(g^2+g)/t\in V$ has constant coefficient $c$. Its square has
only even exponents and is congruent to $c$ modulo $t^2$.
Consequently
\[
 h_2(x)=\left(\frac{g^2+g}{t}\right)^2\in D,
 \qquad h_1(x)=g+t h_2(x)\equiv b\pmod{t^2V}.
\]
Thus $h_1(x)\in D$ as well. Multiplication by $t^2$ proves
\eqref{eq:cusp-explicit-witness}, since the two terms $t^3h_2$
cancel. For example, at $x=t^3$ the values are
\[
 h_2(t^3)=1+t^2+t^6+t^{14},\qquad
 h_1(t^3)=t^3+t^4+t^7+t^{15}.
\]

For the second ring, put $V=\mathbb F_4[[t]]$. Its maximal ideal
is $tV=(t,\alpha t)D$: write every element of $V$ as
$A+\alpha B$ with $A,B\in\mathbb F_2[[t]]\subset D$.
If $x=a+ty\in D$, then $g=f(x)/t=y+ty^2\in V$.
Write $g\equiv b+c\alpha\pmod{tV}$ with $b,c\in\mathbb F_2$.
The relation $\alpha^2+\alpha=1$ gives
\[
 k_2(x)=g^2+g\equiv c\pmod{tV},\qquad
 k_1(x)=g+\alpha k_2(x)\equiv b\pmod{tV}.
\]
Both values therefore belong to $D$, and the last expression for
$k_1$ proves \eqref{eq:residue-extension-explicit-witness}.
All four polynomials have coefficients in the appropriate fraction
field: $t=t^3/t^2$ in the first case and
$\alpha=(\alpha t)/t$ in the second.
\end{proof}

\begin{lemma}[Freeness need not respect polynomial degree]
\label{lem:cusp-no-regular-basis}
Let $D=\mathbb F_2[[t^2,t^3]]$, $V=\mathbb F_2[[t]]$, and
$A=\operatorname{Int}(D)$. Then
\[
 A_{\le2}=D\oplus DX\oplus V(X^2+X),
\]
where $A_{\le2}$ denotes the polynomials of degree at most two.
In particular the ideal of coefficients of $X^2$ in $A_{\le2}$
is $V$, which is not a principal fractional $D$-ideal. Thus $A$
has no $D$-module basis consisting of one polynomial of each
degree, even though $A$ is free as a $D$-module.
\end{lemma}
\begin{proof}
The map $z\mapsto z^2+z$ sends $D$ onto its maximal ideal
$\mathfrak m=t^2V$. Indeed its values lie in $\mathfrak m$, and
for $w\in\mathfrak m$ the convergent series
$z=\sum_{j\ge0}w^{2^j}$ belongs to $\mathfrak m$ and satisfies
$z^2+z=w$.

Every polynomial of degree at most two has a unique expression
$a(X^2+X)+bX+c$ over the fraction field. If it is integer-valued,
evaluation at $0$ and $1$ gives $c,b\in D$. The remaining
condition is $a\mathfrak m\subset D$. For $a\ne0$, write
$a=t^r u$ with $u\in V^\times$. Since $a\mathfrak m=t^{r+2}V$,
this inclusion holds exactly when $r\ge0$: for $r=-1$ or $-2$
the left side contains $t$, and for $r<-2$ it contains an element
of negative valuation. This proves the displayed decomposition.

If $V=aD$, comparison of the least valuations gives $a\in V^\times$.
Multiplying by $a^{-1}$ would then give $D=V$, a contradiction.
Finally, if $(f_n)_{n\ge0}$ were a basis with $\deg(f_n)=n$, the
coefficient of $X^2$ in every polynomial of degree at most two
would be a $D$-multiple of the leading coefficient of $f_2$.
There can be no cancellation with higher basis elements, because
the highest degree in a finite expansion occurs only once.
The ideal $V$ would be principal, contradicting what was proved.
\end{proof}
